\section{Synthesis}
\label{sec:synthesis}

When typing a program using our declarative typing rules, we can
freely apply the subsumption rule to align the (ordered) set of events
performed by the program with a \PAT{} that describes a user's desired
high-level property. Any \emph{synthesis} procedure based on such a
high-level specification must devise a similar ordering alongside the
events in the program it generates. At the same time, each event needs
to align with the specification of its handler in $\Delta$, i.e., its
temporal and data-dependency constraints must be satisfied. Our
solution to this problem is a refinement loop, depicted in
\autoref{fig:intuition-algo}, that iteratively refines the high-level
specification into one that is consistent with these constraints. Each
iteration of this loop targets a single event, adding events before
and after that event so that its dependences are satisfied, i.e., so
that the corresponding handler at that point in the synthesized
program is well-typed. Our loop also employs an regex non-emptiness
check to ensure it represents a generator that produces at least one
feasible trace. After the refinement loop has finished, a
corresponding well-typed generator program can be mechanically
extracted from the refined property.

\subsection{Abstract trace}

Our algorithm normalizes the SRE within violation property and all handler \PAT into a set of \emph{abstract trace}, a sequence of events
$\msg{op}{\phi}$ and kleene stars. This normal form makes it easy to identify
the traces that must precede and follow each event $\msg{op}{\phi}$ in
an SRE's traces.

\begin{definition}[Abstract Trace] An abstract trace $\Pi$ is a SRE without union and top-level Kleene star, defined by the following grammar: %
  {\small
    \begin{alignat*}{2}
      \text{\textbf{Abstract Trace}}& \quad &\Pi  ::= \quad & \epsilon ~|~ se ~|~ \Pi \seqA \Pi ~|~  A^* 
    \end{alignat*}} \noindent
Every SRE can be normalized into a finite set of abstract traces.
\end{definition}

\begin{example}[Abstract trace] The formula encoding violations of a
  Read-Your-Writes policy, \ref{eqn:autoEx}, can be normalized into the
  following abstract trace: %
  {\footnotesize
    \begin{align*}
       \msg{writeRsp}{v =  \Code{x}} \seqA \globalA (\neg
        \msg{writeRsp}{\top}) \seqA   \msg{readRsp}{v =
        \Code{y} \land st = \true \land \Code{y} \not= \Code{x}} \seqA
      \allTL
      \label{eqn:abstraceEx}\tag{$\negSCAPi$}
    \end{align*}}\noindent
  This formula captures the executions of our database example in
  which a successful $\eff{readRsp}$ event carries a value different
  from the last observed $\eff{writeRsp}$ event.
\end{example}


\subsection{Synthesis Algorithm}

\setlength{\textfloatsep}{4pt}
\begin{algorithm}[t!]
  \Params{Ghost variables $\overline{x{:}b}$, handler context
    $\Delta$, and abstract unsafe trace $\Pi$}%
  \Output{Generator $e$, such that
    $\Gamma; \Delta; \emptyset \vdash e:\ha{\epsilon}{\Pi}$} %
  $\Gamma \leftarrow \overline{x{:}\rawnuot{b}{\top}}$; $M \leftarrow \emptyset$
  \tcp*[l]{initialize type context and dependency map}
$(\Gamma, \Pi, M) \leftarrow \S{Refine}(\Gamma, M, \Pi)$ \tcp*[l]{refine events}
$(\Gamma, \Pi) \leftarrow \S{SynLoop}(\Gamma, M, \Pi)$ \tcp*[l]{synthesize loop}
\Return{$\S{DeriveTerm}(\Gamma, M, \Pi)$} \tcp*[l]{derive generator program}
\Procedure{$\mathbf{Refine}(\Gamma, M, \Pi)$}{
  \While(\tcp*[h]{an event that needs to be refined}){exists $\msg{op}{\phi}$ s.t. $\Pi = \Pi_h \seqA \msg{op}{\phi} \seqA \Pi_f$}{
    \If(\tcp*[h]{the events triggered by $\msg{op}{\phi}$ is not examined}){$\eff{op} \not\in \textbf{Dom}(M)$}{
        $(\Gamma, M, \Pi) \leftarrow \forward(\Delta, \Gamma, M, \Pi_h, \msg{op}{\phi}, \Pi_f)$
    }
    \If(\tcp*[h]{the event that triggers $\msg{op}{\phi}$ is not examined}){$\eff{op} \not\in \textbf{Codom}(M) \land \OBS(\eff{op})$}{
        $(\Gamma, M, \Pi) \leftarrow \backward(\Delta,\Gamma,M, \Pi_h, \msg{op}{\phi}, \Pi_f)$
    }
  }
  \Return{$(\Gamma, M, \Pi)$}
}
\caption{Synthesis}
\label{algo:top-syn}
\end{algorithm}

Our top-level synthesis algorithm is shown in \autoref{algo:top-syn}.
Given an (unsafe) abstract trace $\Pi$ and corresponding ghost
variables (e.g., $\Code{x}$ and $\Code{y}$ in \ref{eqn:abstraceEx}) as
input, this nondeterministic algorithm synthesizes a well-typed
$\DSL{}$ generator. The algorithm follows the structure given in
\autoref{fig:intuition-algo}, using a refinement loop (lines $5$ -
$11$) to refine the input abstract trace $\Pi$, synthesize loop (lines $3$), and then
deriving\footnote{The definition of both the SRE normalization
  procedure and $\mathbf{DeriveTerm}$ are provided in the
  \techreport{}.} the final generator from the refined property
(line $10$). Each iteration of this loop nondeterministically chooses
a target event that is used to refine the current abstract trace (line $6$) to preform forward or backward refinement until all events are resolved.
Different choices may result in different event orders, and some of
these choices may cause the algorithm to fail.  Our implementation
resolves this nondeterminism in the algorithm via an efficient
backtracking search procedure that takes the union of all successful
runs in order to capture different orderings.

\paragraph{Event dependencies} The refined abstract trace produced by
our loop must correspond to a well-typed program, i.e., the traces
preceding and following each of its events must be consistent with the
specifications of its corresponding handler. The events that will
precede and follow each event are not known until the loop is finished,
so we cannot simply track the set of observable events, as the
declarative typing rules did via $\Theta$. Instead, each iteration of
the loop detects the unresolved dependencies of a target event in the
current abstract trace and repairs them by inserting appropriate
events before and after it. Intuitively, each observable event
$\eff{op_{child}}$ must follow an operation $\eff{op_{parent}}$, whose handler produces it, forming a tree-like dependency structure
similar to that described by \citet{MessageChain}. The refinement loop
reconstructs these trees from the target node, refining the candidate
abstract trace into one that satisfies these dependencies.

\begin{example}[Event dependency] The refined unsafe abstract trace
  \ref{eqn:autoEx2} contains three distinct sets of events,
  comprised of pairs of requests and corresponding responses: %
  {\footnotesize
    \begin{align*}
      & \msg{writeReq}{v = \Code{x}} \seqA
         \msg{writeReq}{v =\Code{y} \land v \not=\Code{x}} 
         \msg{writeRsp}{v = \Code{y}} \seqA \\
      &\seqA  \msg{readReq}{\top} \seqA 
         \msg{writeRsp}{v = \Code{x}} \seqA
         \msg{readRsp}{v = \Code{y} \land st = \true \land v \not=
        \Code{x}}
    \end{align*}
  }\noindent In contrast, the original trace \ref{eqn:abstraceEx} only
  contains the last two events in \ref{eqn:autoEx2}.
\end{example}

\paragraph{Trace refinement loop} Resolving the dependencies of a
target event requires us to insert events that must precede it and
events that must follow it; the former constitutes the \emph{backward}
synthesis phase of the algorithm, while the latter is subsumed by a
\emph{forward} synthesis pass. The two phases are symmetric: if the
prophecy automaton of the event handler for the operation
$\eff{op_{parent}}$ includes the operation $\eff{op_{child}}$,
performing forward synthesis on $\eff{op_{parent}}$ is the same as
performing backward synthesis on $\eff{op_{child}}$. 
To avoid repeatedly targeting the same event, the algorithm maintains a dependency map $M$ that records the associated child events for each resolved parent event; i.e., $M[\eff{op_{parent}}] = \{\overline{\eff{op_{child}}}\}$.
\footnote{ To enable $M$ to distinguish distinct
  occurrences of events with the same effect operator in the abstract
  trace, we tag each occurence of an operator with a unique
  identifier. For example, \ref{eqn:abstraceEx} with identifiers can
  be$ \msg{writeReq_1}{v = \Code{x}} \seqA \msg{writeReq_2}{v =\Code{y}}...$. } 
The abstract trace corresponds to a well-typed generator program requires all its events are in the domain of $M$, and all \emph{observable} events are in $M$'s codomain, ensuring that no observable event lacks a parent.
% The intersection of these two sets ($\Mfw \cap \Mbw$) contains those events that correspond to well-typed handlers in a generator program.

Depdenecy map $M$ is empty (line $1$) when synthesis starts, and the type context consists
of ghost variables whose qualifiers are $\top$ (line $1$). 
At each iteration of $\mathbf{Refine}$, the algorithm selects a symbolic event $\msg{op}{\phi}$ from the current abstract trace $\Pi$ for refinement. The trace $\Pi$ is partitioned into history ($\Pi_h$) and future ($\Pi_f$) traces surrounding the target event. If the event is not in the domain of $M$, forward synthesis is invoked to introduce necessary child events following it. Conversely, if an observable event is not in the codomain of $M$, backward synthesis is applied to introduce its parent event preceding it. Both synthesis procedures return an updated triple $(\Gamma, M, \Pi)$, reflecting the revised type context, dependency map, and abstract trace. The refinement loop terminates when all events are resolved, i.e., no further events require refinement.

% ; the algorithm terminates once the dependencies of all symbolic events in $\Pi$ are resolved. 
% $\Pi$ is partitioned into the history and future
% traces, $\Pi_f$ and $\Pi_f$, that surround the target event. The
% algorithm tries to perform forward (resp., backward) synthesis on
% these traces, if the event is not in $\Mfw$ (resp., $\Mbw$).  If the
% target operation is generable ($\GEN(\eff{op})$ on line $7$), it is
% the root of a dependency chain, so no additional backward synthesis is
% required and the event is simply added to $\Mbw$ (line $8$). Both the
% forward and backward synthesis routines yield a 6-tuple
% $(\Gamma, \Mfw, \Mbw, \Pi_h, \msgA{op}{\overline{x}}, \Pi_f)$
% that contains updated history and future traces; the refined abstract
% trace at the end of an iteration (line $9$) is simply the
% concatenation of the refined history trace, target event, and refined
% future trace.

\setlength{\textfloatsep}{4pt}
\begin{algorithm}[t!]
\Procedure{$\forward(\Delta, \Gamma, M, \Pi_h, \msg{op}{\phi}, \Pi_f)$}{
    $(\Theta, \overline{z{:}b}\garr\overline{y{:}t}\sarr\rg{\Pi_h'}{\msg{op}{\phi'}}{\Pi_f'}) \in \Delta(eff{op})$\tcp*[l]{select \PAT of $\eff{op}$ from handler context}
    $\Gamma \leftarrow \Gamma, \overline{z{:}\rawnuot{b}{\top}}, \overline{y{:}t}$ \tcp*[l]{add ghost variables and parameters types to type context}
    $\msg{op}{\phi} \leftarrow \msg{op}{\phi \land \phi'}$ \tcp*[l]{merge current regex}
    $\Pi_h \leftarrow \Pi_h \land \Pi_h'$ \tcp*[l]{merge history regex}
    $\Pi_f \leftarrow \Pi_f \land \Pi_f'$ \tcp*[l]{merge prophecy regex}
    \If(\tcp*[h]{non-emptiness check}){$\Gamma \vdash \Pi_h \seqA \msg{op}{\phi} \seqA \Pi_f \not \subseteq \emptyset$}{
        \Return{$(\Gamma, M \cup [\eff{op} \mapsto \Theta], \Pi_h \seqA \msg{op}{\phi} \seqA \Pi_f)$}\;
    }
}
\Procedure{$\backward(\Delta, \Gamma, M, \Pi_h, \msg{op}{\phi}, \Pi_f)$}{
    \tcp{Choose an $\eff{op_{parent}}$ that creates capability $\eff{op}$ and retrieve its \PAT{} from the handler context }
    $(\Theta, \overline{z{:}b}\garr\overline{y{:}t}\sarr\rg{\Pi_h'}{\msg{op_{pa}}{\phi_{pa}}}{\Pi_f'}) \in \Delta(\eff{op_{parent}})$ where $\eff{op} \in \Theta$\;
    $(\Pi_{f_1}' \seqA \msg{op}{\phi'} \seqA \Pi_{f_2}') = \Pi_f'$ \tcp*[l]{The prophecy regex of $\eff{op_{parent}}$ must contains $\eff{op}$}
    $\Gamma \leftarrow \Gamma, \overline{z{:}\rawnuot{b}{\top}}, \overline{y{:}t}$ \tcp*[l]{add ghost variables and parameters types to type context}
    $\msg{op}{\phi} \leftarrow \msg{op}{\phi \land \phi'}$ \tcp*[l]{merge current regex}
    $\Pi_h \leftarrow \Pi_h \land (\Pi_h' \seqA \msg{op_{pa}}{\phi_{pa}} \seqA \Pi_{f_1}')$ \tcp*[l]{merge history regex}
    $\Pi_f \leftarrow \Pi_f \land \Pi_{f_2}'$ \tcp*[l]{merge prophecy regex}
    \If(\tcp*[h]{non-emptiness check}){$\Gamma \vdash \Pi_h \seqA \msg{op_{pa}}{\phi_{pa}} \seqA \Pi_f \not \subseteq \emptyset$}{
        \Return{$(\Gamma, M \cup [\eff{op_{parent}} \mapsto \Theta], \Pi_h \seqA \msg{op_{pa}}{\phi_{pa}} \seqA \Pi_f)$}\;
    }
}
\caption{Forward and Backward Synthesis}
\label{algo:forward-backward}
\end{algorithm}

\paragraph{Forward and backward synthesis} The forward synthesis
subroutine is shown in \autoref{algo:forward-backward}. It first retrieves the
capability and \PAT{} of the target operation $\eff{op}$ from $\Delta$ (line $2$). The algorithm then merges the selected \PAT
into the violation property piecewise. First, the occurence of the
target operation in the current abstract trace $\eff{op}$ is aligned
with its \PAT{} in $\Delta$ (line $4$). Next, the algorithm
merges the history and future traces with the \PAT{}'s history and
prophecy regex (lines $5-6$). In order to guarantee the refined
abstract trace can derive generator program, the algorithm
checks for non-emptiness of the merged abstract trace (line $7$).
Finally, the algorithm returns the refined type context, property, as
well as updated dependency map (line $8$).


The backward synthesis subroutine, shown in \autoref{algo:forward-backward},
is similar to the forward synthesis procedure but works backward from a
target event, insert preceding events into $\Pi_h$ to resolve parent
dependencies. The change in direction results in several differences
with its forward counterpart. The procedure now selects a parent
operator $\eff{op_{parent}}$ whose capability includes the target operator $\eff{op}$ (line $10$), so does its prophecy automata (line $11$). The refined abstract trace needs to align the target operator $\eff{op}$
with its counterpart in the prophecy regex of $\eff{op_{parent}}$:
{\footnotesize
  \begin{align*} \underbrace{ \effL \Pi_h \effR \effL
    \msg{op_{parent}}{\phi_\Code{parent}} \effR \effL \Pi_{f_1}'
    }_\texttt{actual history} \seqA \underbrace{ 
    \msg{op}{\phi'} }_\texttt{actual current } \seqA \underbrace{
    \Pi_{f_2}' \effR }_\texttt{actual prophecy}
  \end{align*}}\noindent 
This is reflected in how the two are merged (line $13$ - $15$). Finally, $\eff{op_{parent}}$ and $\Theta$ are added to the dependency map (line $17$).

\begin{example} We demonstrate the first step of how \ref{eqn:autoEx}
  is refined into \ref{eqn:autoEx2}. The refinement loop begins in the
  following state: %
  { \footnotesize
    \begin{align*}
      \Gamma \equiv\;& \Code{x}\hasoty{int}{\top},
                       \Code{y}\hasoty{int}{\top} \quad \Mfw \equiv \emptyset \quad
                       \Mbw \equiv \emptyset \\
      \Pi \equiv\;& \allTL \seqA 
                    \msg{writeRsp}{v = \Code{x}} \seqA \globalA \neg
                    \msg{writeRsp}{\top} \seqA \msg{readRsp}{v = \Code{y}
                    \land st = \true \land \Code{y} \not= \Code{x}} \seqA \allTL
                    \tag{before iteration $1$}
    \end{align*}}%
  \noindent The first iteration targets the $\eff{readRsp}$ operation.
  Since $\Mfw$ is empty, the algorithm performs forward synthesis on
  $\eff{readRsp}$. No additional events are generated by the handler
  of $\eff{readRsp}$, so no events are added to the abstract
  trace. Since $\eff{readRsp}$ is not generable, the algorithm next
  performs backward synthesis. The signature of $\eff{readReq}$ in
  $\Delta$ uses an intersection \PAT{} whose branches both include
  $\eff{readRsp}$:%
  {\footnotesize
    \begin{align*}
      &\Code{x}\hasnty{int}\garr \rg{\finalA
        \msg{writeReq}{v = \Code{x}} \land
        \neg\nextA\finalA\msg{writeReq}{\top}}{\msg{readReq}{\top}}{\finalA\msg{readRsp}{v
        = \Code{x} \land st = \true}} \tag{$\tau_1$} \\
      & \rg{\neg\finalA
        \msg{writeRsp}{\top}}{\msg{readReq}{\top}}{\finalA\msg{readRsp}{st
        = \false}} \tag{$\tau_2$}
    \end{align*}}\noindent 
  The prophecy automaton of the second branch, $\tau_2$,
  requires $\eff{readRsp}$ to have a \textsf{false} status, which is
  at odds with the current abstract trace. This inconsistency is
  detected by the non-emptiness check, so we backtrack and select the
  next branch, $\tau_1$. This \PAT{} can be merged with the current
  trace, resulting in the following updated values of the refinement
  loop's variables:%
  { \footnotesize
    \begin{align*}
      \Gamma \equiv\;& \Code{x}\hasoty{int}{\top},
                       \Code{y}\hasoty{int}{\top} \quad \Mfw \equiv
                       \{\eff{readRsp}, \eff{readReq}\} \quad \Mbw
                       \equiv \{\eff{readRsp}\} \\
      \Pi \equiv\;& \allTL \seqA \msg{writeReq}{v =\Code{y}}
                    \seqA \globalA \neg \msg{writeReq}{\top} \seqA
                    \msg{readReq}{v = \Code{y}} \seqA \allTL
                    \seqA \msg{writeRsp}{v = \Code{x}} \seqA \\
                     & \globalA \neg \msg{writeRsp}{\top} \seqA \msg{readRsp}{v = \Code{y} \land st = \true \land \Code{y} \not= \Code{x}} \seqA \allTL \tag{after iteration $1$}
    \end{align*}} %
  \noindent
  The refined trace now includes events for $\eff{writeReq}$ and
  $\eff{readReq}$, and the values of the forward and backwards sets
  enable both events to be targeted by the next iteration of the loop.
\end{example}

\subsection{Synthesis Loop}
\label{subsec:synthesis-loop}

\setlength{\textfloatsep}{4pt}
\begin{algorithm}[t!]
\Procedure{$\mathbf{SynLoop}(\Gamma, M, \Pi)$}{
  \ForEach(\tcp*[h]{handle Kleene star from left to right}){$A^*$ s.t. $\Pi = \Pi_h \seqA A^* \seqA \Pi_f$}{
        $\Pi_A \leftarrow \mathbf{InferLoopInv}(\Pi_h, A, \Pi_f)$\;
        $(\Gamma, M, \Pi_h \seqA \textcolor{orange}{\Pi_{\Code{body}}}) \leftarrow \mathbf{Refine}(\Gamma,\Pi_h \seqA \textcolor{orange}{\Pi_A})$ \tcp*[l]{only refine $A$}
        $\Pi \leftarrow \Pi_h \seqA (\Pi_{\Code{body}}^*) \seqA \Pi_f$
  }
  \Return{$(\Gamma,\Pi)$}
}
\caption{Synthesis Loop}
\label{algo:syn-loop}
\end{algorithm}

As illustrated in \autoref{algo:syn-loop}, the algorithm processes all occurrences of $A^*$ from left to right. Each Kleene star $A^*$ in the abstract trace serves as a template for synthesizing loop structures. The subroutine $\mathbf{InferLoopInv}$ infers a suitable loop invariant from $A$ based on the surrounding history and future traces (line $3$). The algorithm then invokes $\mathbf{Refine}$ on the inferred invariant to ensure that the loop body is well-typed (line $4$).

\ZZ{TODO: an example}

\begin{theorem}[Synthesis is Sound]
  \label{theorem:algo-sound}
  The generator synthesized by the algorithm is type-safe with
  respect to our declarative typing rules.
\end{theorem}
